\section{Truth Taught, Power Exercised}

Robert Bellarmine's life acquired coherence from a conviction that truth is a
gift to be learned, ordered, defended, and handed on for salvation.  Humanist
schooling became Jesuit preaching; irregular formation became learned
teaching; the Louvain classroom became the Roman chair of Controversies; large
volumes became short catechisms; curial office was interrupted by residential
pastoral rule; late controversy returned to retreat and preparation for death.
The sequence shows unusual continuity between intellectual and spiritual
work.

It also shows why learning cannot be separated from the power that receives
it.  Bellarmine's distinctions could limit a claim of direct papal temporal
dominion, counsel against defining a disputed theory of grace, and leave open
the reinterpretation of Scripture before a true demonstration.  The same
institutional confidence supported censorship, intervention against rulers,
death for some heretics, a consequential role in Bruno's process, and the 1616
restriction imposed on Galileo.  Nuance is not exoneration; precision is not
indictment by association.

Capua prevents the curial record from becoming the whole man.  Bellarmine
accepted a residential see, convened synods and a provincial council, visited,
preached, and described care for the poor.  The catechisms and retreat books
show why doctrine mattered to him beyond victory: it was meant to form persons
for worship, moral action, and eternal life.  The source limits prevent those
pastoral goods from becoming a defense brief.  Much of their detail comes
through his own retrospective account and later official praise, while a full
Capuan archive remains to be collated.

His sanctity, as received by the Catholic Church, is therefore not a claim that
history has no remainder.  A Doctor can teach eminently without every
political judgment being permanent; a saint can inhabit institutions whose
coercion must be named; an opponent can suffer injustice without every later
story about the case becoming true.  Bellarmine's own best habits---exact
locus, distinction of powers, attention to genre, correction of printed
work---supply tools for reading him responsibly.

The unresolved questions remain material.  The complete Bruno propositions
and Roman dossier do not survive; the 1616 precept record remains disputed;
the Sixtine Index procedure, Vulgate meetings, \emph{de auxiliis} memorials,
Capuan acts, conclave ballots, reciprocal political dossiers, early process
testimony, and relic history need fuller collation.  Responsible memory does
not fill those spaces with innocence or guilt invented for dramatic symmetry.

Bellarmine's durable significance lies in the union that makes him difficult:
he taught through institutions because he believed truth should order a
Christian people.  That union made possible remarkable catechetical,
theological, pastoral, and spiritual service.  It also gave disputed judgments
real coercive reach.  A source-first life receives the whole inheritance by
honoring the saint, listening to opponents and victims, and refusing to make
either reverence or criticism an excuse for imprecision.
